\section{The thirteen-point case and the direct large range}
\label{sec:n13-infinity}

Write
\[
  F(n)=1+\binom{n-1}{2}-\floor{\frac{n-1}{2}}.
\]
We finish the finite range at $n=13$ and then prove the lower bound
$C\geq F(n)$ directly for every $n\geq14$.  The large-range argument is
not recursive: it uses neither a value at a preceding integer nor a
distinguished base value.

For a point set $P$, let $C_s$ and $L_s$ be the numbers of maximal proper
circle blocks and maximal line blocks of size $s$, respectively, and put
$B_s=C_s+L_s$ and $C=\sum_s C_s$.  Every triple has a unique maximal
generalized-block owner.  We shall repeatedly use
\begin{align}
 \mathsf T&:=\sum_{s\geq3}\binom{s}{3}B_s=\binom n3,
                                                        \label{eq:large-T}\\
 \mathsf P&:=\sum_{s\geq3}s(4-s)B_s\geq3n,             \label{eq:large-P}\\
 \mathsf D&:=\sum_{s\geq3}s(s-4)C_s
       +\sum_{s\geq3}2s(s-2)L_s\leq n(n-4).           \label{eq:large-D}
\end{align}
Indeed, inversion at a selected point followed by Melchior's inequality
gives
\[
 \sigma_p=d_3(p)-3-\sum_{s\geq5}(s-4)d_s(p)\geq0.
\]
Summing this inequality gives \eqref{eq:large-P}.  Applying Melchior once
more after restoring the centre of inversion gives
\(\sum_s s\ell_s(p)\leq n-1+\sigma_p\); summing over $p$ and eliminating
the $\sigma_p$ gives \eqref{eq:large-D}.  These are geometric consequences
of the real projective-plane inequality of Melchior~\cite{Melchior1940},
not merely formal identities.

\subsection{The complete thirteen-point dispatch}

Suppose, towards a contradiction, that $|P|=13$ and $C\leq60$.  Let $m$
be the largest size of a proper circle block.

\begin{lemma}[Outer routing at thirteen points]
\label{lem:n13-outer}
Under the preceding assumptions, $m\in\{5,6\}$.  Moreover every line block
has size at most six.
\end{lemma}

\begin{proof}
If a line has $s$ selected points and $q=13-s$, the rich-line pencil gives
\begin{equation}
 q\binom{s}{2}-\binom q2\floor{\frac s2}.           \label{eq:n13-line-pencil}
\end{equation}
For $7\leq s\leq12$ its minimum is $66$.  Thus no such line occurs, while
$s=13$ is excluded by noncollinearity.
For a proper $m$-circle, the corresponding Bonferroni pencil is
\begin{equation}
 1+q\left(\binom m2-\floor{\frac m2}\right)
   -\binom q2\floor{\frac m2},\quad q=13-m.          \label{eq:n13-circle-pencil}
\end{equation}
For $m=7,8,9,10,11,12$ this equals, respectively,
\[
 64,81,105,106,96,61,
\]
so $m\geq7$ is impossible.  The case $m=13$ is excluded by the standing
non-concyclicity hypothesis.

It remains to exclude $m\leq4$.  Besides \eqref{eq:large-T} and
\eqref{eq:large-P}, global Melchior for line blocks gives
\[
 \mathsf G:=3L_3+\sum_{s\geq4}
 \left(\binom s2+s-3\right)L_s\leq\binom{13}{2}-3.
\]
Since line sizes are at most six, the fixed combination
\[
 \frac14\mathsf T+\frac7{60}\mathsf P-\frac15\mathsf G
\]
has coefficients
\[
\begin{array}{c|cc|rrrr}
 &C_3&C_4&L_3&L_4&L_5&L_6\\ \hline
 &3/5&1&0&-2/5&-29/60&0.
\end{array}
\]
It is therefore at most $C$, whereas its numerical lower bound is
\[
 \frac14\binom{13}{3}+\frac7{60}(3\cdot13)
 -\frac15\left(\binom{13}{2}-3\right)
 =\frac{1221}{20}>60.
\]
Thus $m\leq4$ is impossible as well.
\end{proof}

\subsubsection{The six-circle branch}

Fix a six-point circle $\Gamma$ and let $X=P\setminus\Gamma$, so
$|X|=7$.  If $e=xy\in\binom X2$, let $w_e$ be the number of proper
circles through $x,y$ and two points of $\Gamma$, and put
\begin{equation}
 W=\sum_{e\in\binom X2}w_e
   =\sum_j\binom j2c(2,j).                          \label{eq:n13-W}
\end{equation}
Here $c(i,j)$ counts nonselected circle blocks meeting $\Gamma$ in $i$
points and $X$ in $j$ points; $l(i,j)$ denotes the analogous number of
line blocks.
For either class $Y\in\{\Gamma,X\}$, define its pivot row by
\begin{equation}
 M_Y=\sum_B |B\cap Y|(4-|B|),                       \label{eq:n13-class-pivot}
\end{equation}
where the sum is over all maximal blocks, including $\Gamma$.
Summing the pointwise Melchior slack over $p\in Y$ gives
\[
 M_Y=3|Y|+\sum_{p\in Y}\sigma_p.
\]
Consequently
\begin{equation}
 M_\Gamma\geq18,\quad M_X\geq21.                    \label{eq:n13-class-pivot-bounds}
\end{equation}

\begin{lemma}[The opening inequality]
\label{lem:n13-gamma6-opening}
If $C\leq60$, then $W\geq48$.
\end{lemma}

\begin{proof}
For thirteen points, \eqref{eq:large-D} reads $\mathsf D\leq117$.
Let $O=3C_3$ be ordinary-circle incidence.  For every pivot $p$, the
inverted arrangement in \cref{cor:pivot-lenchner} has
$q=13-1=12\neq6$ lines.
Thus at least three ordinary proper circles pass through each pivot.
Summing over the thirteen pivots, and using
\eqref{eq:n13-class-pivot-bounds}, gives
\[
 O\geq39,\quad M_X\geq21.
\]
Weight the four triple partitions, in order of their number of
$\Gamma$-points, by
\[
 (141,340,539,141).
\]
Add $79O+199M_X$ and subtract $1344C+414W$.  For a nonselected block
write $g=|B\cap\Gamma|$, $x=|B\cap X|$, and $s=g+x$.  Its coefficient is
\begin{align*}
 A(g,x)={}&141\binom x3+340g\binom x2
  +539\binom g2x+199x(4-s)\\
 &+237\mathbf{1}_{\{B\text{ is a }3\text{-circle}\}}
  -1344\mathbf{1}_{\{B\text{ is a circle}\}}\\
 &-414\mathbf{1}_{\{B\text{ is a circle and }g=2\}}\binom x2.
\end{align*}
Direct factorisation, recorded in
\cref{subsec:n13large-certificates}, gives $A(g,x)\leq123D_B$ for every
geometrically possible block, including equality at the selected circle.
Since the four triple totals are $35,126,105,20$, summing gives
\begin{equation}
 123\mathsf D\geq33810-414W.                       \label{eq:n13-opening-dual}
\end{equation}
Together with $\mathsf D\leq117$, this first yields $W\geq47$.

At $W=47$, retaining every nonnegative slack in the same calculation gives
\[
 R+79(O-39)+199(M_X-21)+1344(60-C)\leq39,
\]
where every positive local residual contributing to $R$ is at least $96$.
Thus all four terms vanish.  For the zero-residual circle types
\[
 c(0,3),c(1,2),c(1,5),c(2,1),c(2,2),c(2,3),
\]
write their multiplicities, in order, as $(a,b,q,d,e,f)$.  The three
zero-residual line types and their multiplicities are
\[
\begin{array}{c|ccc}
\text{type}&l(0,3)&l(1,2)&l(2,1)\\ \hline
\text{multiplicity}&p_0&p_1&p_2.
\end{array}
\]
The circle count, ordinary-circle count, $W$, and the three triple
partitions reduce to
\begin{gather*}
 a+b+d=13,\quad q+e+f=46,\quad e+3f=47,\quad q=2f-1,\\
 p_0=45-a-21f,\quad p_1=42-b-20f,\quad p_2=11-d+3f.
\end{gather*}
Consequently
\[
 M_X=3a+2b+d+3p_0+2p_1+p_2-10q-3f=240-123f.
\]
The equality $M_X=21$ would force $f=73/41$, a contradiction.
\end{proof}

We now invoke the six-marked-conic interface
\cref{lem:six-conic-events}.  In the present notation, its first clause
gives $w_e\leq3$ and at most four projective involution signatures for the
\emph{full} edges $w_e=3$.  Each signature support is a matching in $X$,
and at most three full-signature centres lie on one trace line.  A pair of
support edges of the same signature is a \emph{repetition event}.
The trace dictionary \cref{lem:circle-lift-trace} and the second clause of
\cref{lem:six-conic-events} give every event a unique maximal
line-or-circle host.  If that host has $b$ outsider points, its event
capacity is
\begin{equation}
 k_0(b)=\min\left\{2\binom b4,
       3\binom{\floor{b/2}}2\right\};\quad
 k_0(4),k_0(5),k_0(6)=2,3,9.                        \label{eq:n13-host-cap}
\end{equation}
This is exactly the capacity \eqref{eq:conic-event-cap}.  The host cannot be
a line meeting $\Gamma$ or a circle meeting $\Gamma$ once.  A host circle
meeting $\Gamma$ in two points is admissible only for the signature
containing that marked pair.  These assertions follow by uniqueness of the
lift plane in \cref{lem:circle-lift-trace} and by the fact that two proper
circles sharing three points coincide.

\begin{lemma}[Capacity-gap lemma]
\label{lem:n13-capacity-gap}
Put $r=W-48$.  Then $0\leq r\leq6$.  If $R_{\rm rep}$ is the number of
repetition events, there is a nonnegative integer
\[
 E=2h+3k+9A_6+2\ell+2g,
\quad
 (h,k,A_6,\ell,g)=(c(0,4),c(0,5),c(0,6),l(0,4),c(2,4)),
\]
such that
\begin{equation}
 R_{\rm rep}\leq E,\quad 4E\leq6r+7+8A_6,\quad A_6=0.
                                                               \label{eq:n13-capacity-gap}
\end{equation}
\end{lemma}

\begin{proof}
By \cref{lem:n13-gamma6-opening}, $W\geq48$.  If $F_3$ denotes the number
of full edges, then
\[
 W\leq2\binom72+F_3.
\]
The four signature supports are matchings of size at most three, so
$F_3\leq12$ and hence $W\leq54$.  This proves $0\leq r\leq6$.

Two possible host types must first be removed explicitly.  Write
$\lambda_5=l(0,5)$ and $\lambda_6=l(0,6)$.  On a line of type $l(0,x)$,
the second-dual residual is $x(x-3)(14-x)$, hence it is $90$ at $x=5$
and $144$ at $x=6$.  Before forced-zero variables are suppressed, the
full second added-centre slack inequality therefore contains
\[
 90\lambda_5+144\lambda_6+\text{other nonnegative terms}
 \leq12+6r\leq48.
\]
Hence $\lambda_5=\lambda_6=0$.  Host uniqueness and admissibility now leave
only
\[
 c(0,4),\ c(0,5),\ c(0,6),\ l(0,4),\ c(2,4).
\]
Their capacities from \eqref{eq:n13-host-cap} are respectively
$2,3,9,2,2$, so this exhaustive list gives $R_{\rm rep}\leq E$.

For the numerical inequality, introduce
\[
 \begin{gathered}
 y=c(1,3),\ z=c(1,4),\ f=c(2,3),\quad
 w=l(1,3),\ t=l(2,2),\ p=l(2,3),\\
 d=60-C,\quad j=C_3-13,\quad
 u=M_X-21,\quad v=M_\Gamma-18,
 \end{gathered}
\]
all nonnegative integers.  The second added-centre calculation has the
exact retained-slack budget
\begin{align}
 &12h+26k+36A_6+3y+5z+2f+40\ell+31w+22t+48p
       +6u+3v+36d\leq12+6r.                       \label{eq:n13-second-budget}
\end{align}
Let $N=j+y+z$, $a=y+z$, and define
\begin{align*}
 K_0={}&N+2u+v+4h+8k+12A_6+10\ell+8d+4t+7w,\\
 Q_0={}&N+u+4h+8k+12A_6+8\ell+8d+2t+5w+7p.
\end{align*}
The eight conservation rows displayed in
\cref{subsec:n13large-certificates} give an integer $\mu\geq0$ and
$\delta=7\mu-4j$ satisfying
\begin{align}
 K_0+7j-2r&=14\mu,& Q_0+4\mu-2j&\leq3+2r,
                                                               \label{eq:n13-mu}\\
 3\delta&\leq6-k-5\ell-6d-5t-5w-24p.             \label{eq:n13-delta}
\end{align}
They also give, with
\[
 J=21+18r-(24g+24h+36k+84A_6+24\ell),
\]
the exact identity
\begin{align}
7J={}&49f+12j-60\delta+14a+7y+63u+14v+105k
 +168\ell+252d-161p+42t+189w.                    \label{eq:n13-J}
\end{align}
 Its right side is nonnegative.  If $p>0$, \eqref{eq:n13-delta} gives
\[
 \delta\leq2-8p,\quad
 -60\delta-161p\geq319p-120\geq199>0.
\]
Thus the only apparently negative pair of terms is already positive before
the remaining nonnegative terms are restored.
If $p=0$ and $\delta\leq0$, this is immediate.  For $\delta=1$, the
congruence $\delta=7\mu-4j$ gives $j\equiv5\pmod7$, hence
$12j-60\geq0$.  For $\delta=2$, \eqref{eq:n13-delta} forces
$k=\ell=d=t=w=0$ and $j\equiv3\pmod7$.  If $j\geq10$, then
$12j-120\geq0$.  For $j=3$ one has $\mu=2$, and
\eqref{eq:n13-mu} gives
\[
 a+u+4h+12A_6\leq2r-2.
\]
The two eliminated conservation identities then read $v\geq6-u$ and
\[
 f=(y-a)-2\bigl(v-(6-u)\bigr).
\]
Since $f\geq0$ and $y\leq a$, equality holds throughout:
$f=0$, $y=a$, and $v=6-u$.  The right side of
\eqref{eq:n13-J} is therefore $21a+49u\geq0$.  Hence $J\geq0$ in every
case.

Since
\[
 12E-24A_6=24g+24h+36k+84A_6+24\ell,
\]
$J\geq0$ is precisely the middle inequality of
\eqref{eq:n13-capacity-gap}.

It remains to remove $A_6$.  A first fixed added-centre row gives
\[
 2388A_6\leq453+414r,
\]
and hence $A_6=0$ for $0\leq r\leq4$.  For $r=5,6$, the number of full
edges is at least $6+r\geq11$.  Four matching supports, each of size at
most three, must therefore all have size at least two.  If a six-outsider
circle existed, every other event host would share at least three of its
six outsider points and hence would coincide with it.  All four signature
centres would then lie on its trace line, contradicting the three-centre
cap.  Thus $A_6=0$ also for $r=5,6$.
\end{proof}

The number $F_3$ of full outsider edges satisfies
\(W\leq2\binom72+F_3\), so $F_3\geq6+r$.  Balancing these edges among four
matching supports of capacity three gives
\begin{equation}
 R_{\rm rep}\geq\rho(r),\quad
\begin{array}{c|rrrrrrr}
r&0&1&2&3&4&5&6\\ \hline
\rho(r)&2&3&4&6&8&10&12.
\end{array}                                        \label{eq:n13-rho}
\end{equation}
Together, \eqref{eq:n13-capacity-gap} and \eqref{eq:n13-rho} immediately
exclude $r=0,4,5,6$.  At $r=1$, \eqref{eq:n13-second-budget} forces
$k=A_6=\ell=0$, so $E=2(h+g)$ is even; the capacity gap gives $E\leq3$,
and hence $E\leq2<\rho(1)$.  Only $r=2,3$ remain.

\begin{lemma}[Three-residue terminal]
\label{lem:n13-three-residue}
The equality layers $r=2$ and $r=3$ are empty.
\end{lemma}

\begin{proof}
For $r=2$, \eqref{eq:n13-rho} and the capacity gap give
\[
 4\leq E\leq\floor{\frac{19}{4}}=4;
\]
for $r=3$ they give
\[
 6\leq E\leq\floor{\frac{25}{4}}=6.
\]
Thus $E=4,6$, respectively.  Since $A_6=0$, the definition of $J$ becomes
$J=21+18r-12E$, and hence $J=9$ for $r=2$ and $J=3$ for $r=3$.

We next justify the reduction of variables.  At $r=2$, the right side of
\eqref{eq:n13-second-budget} is $24$; its coefficients force
\[
 k=\ell=d=w=p=0.
\]
At $r=3$, its right side is $30$, so
$\ell=d=w=p=0$ and $k\leq1$.  But here
\[
 E-3k=2(h+g)
\]
is even, whereas $E=6$; hence $k$ cannot equal one and again $k=0$.
Together with $A_6=0$, this leaves precisely the nonnegative integers
$a,y,u,v,h,t,f,j,\mu$, with $0\leq y\leq a$.  They satisfy
\begin{align}
 \delta&=7\mu-4j,                                  \label{eq:n13-res-delta}\\
 v&=2\delta+2r-a-2u-4h-4t,\nonumber\\
 f&=2j+\mu+y-v+4h-5t-2r,                           \label{eq:n13-res-vf}\\
 a+u+4h+2t+4\mu-j&\leq3+2r,                       \label{eq:n13-res-budget}\\
 J&=7a+48h+90j-123\mu-24r-9t+19u+8y.             \label{eq:n13-res-J}
\end{align}
The capacity inequality gives $\delta\leq2$.  Negative $\delta$ is
impossible: for $r=3$, the term $-60\delta$ in the uneliminated identity
\eqref{eq:n13-J} already exceeds $7J=21$; for $r=2$, the value
$\delta=-1$ forces $j\equiv2\pmod7$, and then
$-60\delta+12j\geq84>63=7J$.  More negative values only increase this
lower bound.  Thus $\delta\in\{0,1,2\}$.

Moreover,
\[
 4\mu-j=\frac{4\delta+9j}{7}.
\]
Dropping the other nonnegative terms in \eqref{eq:n13-res-budget} gives
$4\delta+9j\leq49$ for $r=2$ and at most $63$ for $r=3$.  The congruence
$j\equiv5\delta\pmod7$ therefore leaves
\begin{equation}
 (\delta,j,\mu)\in\{(0,0,0),(1,5,3),(2,3,2)\}.    \label{eq:n13-three-residues}
\end{equation}
For completeness, the sole apparent extra triple is $(0,7,4)$ when
$r=3$, but then the nonnegative right side of \eqref{eq:n13-J} contains
$12j=84>21$.

For $(0,0,0)$, the two inequalities $v,f\geq0$ give
\[
 a+2u+4h+4t\leq2r,\quad
 a+y+2u+8h-t\geq4r.
\]
Since $y\leq a$, subtraction yields
\[
 2r\leq y+4h-5t\leq a+4h-5t\leq2r-2u-9t.
\]
Hence $u=t=0$, $y=a=2r-4h$, and
$J=6r-12h$.  The pairs $(r,J)=(2,9),(3,3)$ would respectively make
$h=1/4,5/4$.

For $(1,5,3)$, the two relevant rows become
\begin{align*}
 a+u+4h+2t&\leq2r-4,\\
 J+24r-81&=7a+48h-9t+19u+8y.
\end{align*}
When $r=2$, the budget makes every variable zero, whereas the left side
of the identity is $-24$.  When $r=3$, the budget is two: for $t=0$ the
right side is nonnegative, and for $t=1$ it equals $-9$, while the left
side equals $-6$.

Finally, for $(2,3,2)$,
\begin{align*}
 a+u+4h+2t&\leq2r-2,\\
 J+24r-24&=7a+48h-9t+19u+8y. 
\end{align*}
For $r=2$ the left side is $33$ and the budget is two.  Thus $h=0$.
If $t=1$, the right side is $-9$; if $t=0$, parity forces
$a+u=1$, and the right side is at most $19$.

For $r=3$ the left side is $51$ and the budget is four.  The value $h=1$
gives $48$, while $t=2$ gives $-18$ and $t=1$ gives at most $29$.
Thus $h=t=0$.  Parity leaves $a+u\in\{1,3\}$.  The first value gives at
most $19$; for the second, substituting $u=3-a$ gives
$4(2y-3a)=-6$.  This is impossible.
\end{proof}

\begin{proposition}
\label{prop:n13-gamma6}
If a thirteen-point set has a largest proper circle of size six, then
$C\geq61$.
\end{proposition}

\begin{proof}
Assume $C\leq60$.  The capacity-gap lemma places
$r=W-48$ in $\{0,\ldots,6\}$ and excludes $r=0,1,4,5,6$, while
\cref{lem:n13-three-residue} excludes $r=2,3$.
\end{proof}

It remains to close the five-circle value left by \cref{lem:n13-outer}.

\begin{lemma}[The five-circle transfer]
\label{lem:n13-gamma5}
If a thirteen-point set has no proper circle larger than five, then
$C\geq61$.
\end{lemma}

\begin{proof}
Assume $C\leq60$.  By \eqref{eq:n13-line-pencil}, every line has size at
most six, so only $B_3,B_4,B_5,B_6$ occur and $B_6=L_6$.  Set
\begin{align*}
 K&=\sum_{p\in P}\sigma_p,\quad
 Q=3B_3-5B_5-12B_6=39+K,\\
 S&=3L_3+4L_4+10L_5+12L_6,\\
 A&=5L_3+12L_4+20L_5+28L_6.
\end{align*}
The triple partition and the two applications of Melchior described above
give
\begin{align}
 B_3+4B_4+10B_5+20B_6&=286,                       \label{eq:g5-triples}\\
 12B_4+35B_5+72B_6&=819-K,                         \label{eq:g5-block}\\
 A&\leq4C-169+C_5,\quad A\geq\frac53S,             \label{eq:g5-A}\\
 24\cdot286+27Q&\leq35(S+3C).                     \label{eq:g5-transfer}
\end{align}
The last inequality is coefficientwise:
\[
\begin{array}{c|rrrr}
s&3&4&5&6\\ \hline
\text{left side}&105&96&105&156\\
\text{right side, circle}&105&105&105&--\\
\text{right side, line}&105&140&350&420.
\end{array}
\]
The dash records $C_6=0$ in this branch.  Thus, on every type that can
occur, no separation or optimisation theorem is hidden in
\eqref{eq:g5-transfer}.

From \eqref{eq:g5-block}, $C_5\leq B_5\leq23$.  Equations
\eqref{eq:g5-A}--\eqref{eq:g5-transfer} now give
\[
 377-5C\leq A\leq4C-146,
\]
and hence $C\geq523/9$, so only $C=59,60$ remain.

For these two endpoints let $L=\sum_{s=3}^6L_s$.  Exact elimination from
\eqref{eq:g5-triples}--\eqref{eq:g5-A} gives
\begin{align}
 4L&=325-4C-B_5-4B_6+K,                            \label{eq:g5-L}\\
 9B_5&\geq2301-36C+5K+72B_6.                      \label{eq:g5-B5}
\end{align}
Combining \eqref{eq:g5-block}, $B_4\geq0$, and \eqref{eq:g5-B5} yields
\[
 \begin{array}{ll}
 C=59:&23K+396B_6\leq147,\\
 C=60:&46K+792B_6\leq609.
 \end{array}
\]
Thus $B_6=0$.  Reducing \eqref{eq:g5-block} modulo $12$, then using
\eqref{eq:g5-B5} and $B_4\geq0$, gives exactly
\begin{equation}
 (B_3,B_4,B_5,B_6)=(48+2K, 7-3K, 21+K, 0),
 \quad K\in\{0,1,2\}.                              \label{eq:g5-three-endpoints}
\end{equation}
At every point $p$, pair partition and the definition of $\sigma_p$ give
\begin{equation}
 7d_5(p)+3d_4(p)=63-\sigma_p.                      \label{eq:g5-point}
\end{equation}

If $K=0$, \eqref{eq:g5-point} and $d_4(p)\leq B_4=7$ force
$d_4(p)\in\{0,7\}$.  Since $\sum_p d_4(p)=4B_4=28$, four points lie in
all seven four-blocks, making two distinct maximal blocks share four
points.  If $K=1$, twelve zero-slack points have $(d_5,d_4)=(9,0)$ and
the remaining point has $(8,2)$; hence $\sum_p d_5(p)=116$, rather than
$5B_5=110$.  Finally, if $K=2$, any point with positive slack must have
$d_4(p)\geq2$, while \eqref{eq:g5-three-endpoints} gives $B_4=1$.
All three cases are impossible.
\end{proof}

\begin{theorem}\label{thm:n13-value}
The thirteen-point value is
\[
 f(13)=61.
\]
\end{theorem}

\begin{proof}
For a hypothetical set with $C\leq60$, \cref{lem:n13-outer} leaves only
$m=5$ and $m=6$.  These are excluded by \cref{lem:n13-gamma5} and
\cref{prop:n13-gamma6}, respectively.  Hence $C\geq61=F(13)$.  The
circle-with-centre construction from the foundational section has exactly
$F(13)$ proper circles, proving equality.
\end{proof}

\subsection{A uniform half-cap and one master for
  \texorpdfstring{$n\geq15$}{n >= 15}}

\begin{lemma}[Uniform half-cap]
\label{lem:large-half-cap}
Let $n\geq15$ and suppose $C\leq F(n)-1$.  Every maximal line or proper
circle block then has size at most
\[
 H=\floor{\frac n2}.
\]
\end{lemma}

\begin{proof}
Let a block have size $s$, put $q=n-s\geq1$, and write
$h=\floor{s/2}$.  The rich-circle and rich-line pencils are
\begin{align}
 R(n,s)&=1+q\left(\binom s2-h\right)-\binom q2h,   \label{eq:large-rich-circle}\\
 Q(n,s)&=q\binom s2-\binom q2h=R(n,s)-1+qh.        \label{eq:large-rich-line}
\end{align}
We show $R(n,s)\geq F(n)$ whenever $s\geq H+1$.  The following four rows
cover all parity choices; in each row $c\geq0$.
\[
\begin{array}{c|c|c|l}
s&q&\text{parameter relation}&R(n,s)-F(n)\\ \hline
2a&2b&a=b+1+c&
(4b-2)c^2+(6b^2-3b)c+2b^3-5b^2+b\\
2a&2b+1&a=b+1+c&
b\{4c^2+(6b-1)c+2b^2-5b-3\}\\
2a+1&2b&a=b+c&
(4b-2)c^2+(6b^2-7b+2)c+2b^3-7b^2+4b\\
2a+1&2b+1&a=b+1+c&
b\{4c^2+(6b+3)c+2b^2-b-1\}.
\end{array}                                                    \tag{*}
\]
The condition $n\geq15$ is, row by row,
\(2b+c\geq7,6,7,6\).  For $b\geq3$ all relevant coefficients in the
first three rows are nonnegative, and the fourth row is nonnegative for
every $b\geq1$.  The remaining possibilities reduce to
\[
\begin{array}{c|cc}
&b=1&b=2\\ \hline
\text{row 1}&2c^2+3c-2\ (c\geq5)&6c^2+18c-2\ (c\geq3)\\
\text{row 2}&4c^2+5c-6\ (c\geq4)&2(4c^2+11c-5)\ (c\geq2)\\
\text{row 3}&2c^2+c-1\ (c\geq5)&6c^2+12c-4\ (c\geq3),
\end{array}
\]
all strictly positive.  In the two odd-$q$ rows, $b=0$ gives equality
$R=F$.  Thus a circle of size $s\geq H+1$ already determines at least
$F(n)$ circles.  Equation \eqref{eq:large-rich-line} and $qh\geq1$ give
the same conclusion for a line.  Either conclusion contradicts
$C\leq F(n)-1$.
\end{proof}

Under the cap $s\leq M$, inversion leaves at most $M-1$ collinear images.
For $M=\floor{n/2}$,
\[
 M-1\leq\frac{2(n-1)}3.
\]
The Langer--de Zeeuw inequality \cref{thm:langer}, applied after each
inversion and summed over all centres, therefore gives
\begin{equation}
 \mathsf L:=\sum_{s\geq3}s(s-1)B_s
 \geq\frac{n(n-1)(n+2)}3.                          \label{eq:large-L}
\end{equation}

\begin{lemma}[The $S_M$ master]
\label{lem:large-SM}
Assume every maximal block has size at most $M\geq3$ and that
\eqref{eq:large-L} holds.  Then
\begin{equation}
 C\geq A(n,M):=
 \frac{\binom n3}{2M}+\frac{n(M+3)}{6M}
 +\frac{(M-2)n(n-1)(n+2)}{36M}-\frac{n(n-4)}9.
                                                               \label{eq:large-A}
\end{equation}
\end{lemma}

\begin{proof}
Consider the single functional
\begin{equation}
 \mathsf S_M=\frac{\mathsf T}{2M}
 +\frac{M+3}{18M}\mathsf P
 +\frac{M-2}{12M}\mathsf L-\frac19\mathsf D.      \label{eq:large-SM}
\end{equation}
If $k_C(s)$ and $k_L(s)$ are its coefficients on an $s$-circle and an
$s$-line, respectively, direct factorisation gives
\begin{align}
 1-k_C(s)&=\frac{(M-s)(s-4)(s-3)}{12M},            \label{eq:large-circle-slack}\\
 k_L(s)&=\frac{s(s-3)(-7M+3s-12)}{36M}.            \label{eq:large-line-slack}
\end{align}
Thus $k_C(3)=k_C(4)=1$, $k_C(s)\leq1$ for $5\leq s\leq M$, and
$k_L(s)\leq0$ for every $3\leq s\leq M$.  Hence
$\mathsf S_M\leq C$.  Substitution of
\eqref{eq:large-T}, \eqref{eq:large-P}, \eqref{eq:large-D}, and
\eqref{eq:large-L}, with attention to the negative coefficient of
$\mathsf D$, gives \eqref{eq:large-A}.
\end{proof}

Taking $M=\floor{n/2}$ in \eqref{eq:large-A} yields
\begin{align}
 A(2r,r)-F(2r)&=\frac{(2r-1)(r^2-9r+13)}9,
                                                               \label{eq:large-even-gap}\\
 A(2r+1,r)-F(2r+1)&=
 \frac{4r^4-32r^3+37r^2+9}{18r}.                  \label{eq:large-odd-gap}
\end{align}
The first is positive for $r\geq8$; after $r=8+t$ its quadratic factor is
$t^2+7t+5$.  In the second, after $r=7+t$, the numerator is
\[
 4t^4+80t^3+541t^2+1302t+450,
\]
so it is positive for $r\geq7$.  Together with
\cref{lem:large-half-cap,lem:large-SM}, this proves $C\geq F(n)$ for
every $n\geq15$.

\subsection{The exact dichotomy at fourteen points}

\begin{proposition}[Seven-circle dichotomy]
\label{prop:n14-dichotomy}
Every admissible fourteen-point set determines at least $73$ proper
circles.
\end{proposition}

\begin{proof}
Assume $C\leq72$.  The pencil bounds
\eqref{eq:large-rich-circle}--\eqref{eq:large-rich-line} show that proper
circle blocks have size at most seven and line blocks have size at most
six.  If $C_7=0$, all blocks have size at most six.  Every inverted line
then contains at most five of the thirteen image points, and
$5\leq2\cdot13/3$, so \cref{thm:langer} supplies \eqref{eq:large-L}.
The $S_M$ master with $M=6$ gives
\[
 C\geq A(14,6)=\frac{3899}{54}>72.
\]

It remains to consider a selected seven-circle $\Gamma$ with seven-point
complement $X$.  For every block put
$g=|B\cap\Gamma|$, $x=|B\cap X|$, and $s=g+x$, and define
\begin{equation}
 W=\sum_{\substack{B\text{ a proper circle}\\g=2}}\binom x2.
                                                               \label{eq:n14-W}
\end{equation}
The four triple totals according to their number of $\Gamma$-points are
$\binom7i\binom7{3-i}$.  Put
\[
 M_\Gamma=\sum_B g(4-s),\quad M_X=\sum_B x(4-s),
\]
where the sums include the selected block.  Summed pivot Melchior over each
of the two seven-point classes gives $M_\Gamma,M_X\geq21$, while
\eqref{eq:large-D} gives $\mathsf D\leq140$.

For every block set
\begin{align*}
 A_B={}&6\binom x3+9g\binom x2+12\binom g2x+6\binom g3
 +(3g+6x)(4-s)\\
 &-36\mathbf{1}_{\{B\text{ is a circle}\}}
 -6\mathbf{1}_{\{B\text{ is a circle and }g=2\}}\binom x2.
\end{align*}
For a nonselected block one has $g\leq2$, and the residual $4D_B-A_B$
factors as
\begin{equation}
\begin{array}{c|ccc}
&g=0&g=1&g=2\\ \hline
\text{circle}&(x-3)(-x^2+10x-12)&
 (5-x)(x-2)(2x-3)/2&(4-x)(x-2)(x-1)\\
\text{line}&x(14-x)(x-3)&
 (x-2)(-2x^2+21x+17)/2&(x-1)(-x^2+7x+12).
\end{array}                                                     \label{eq:n14-residuals}
\end{equation}
Under the block caps every line residual is nonnegative; a negative circle
residual has magnitude at most $3\binom x3$.  Outsider triples have unique
owners, so the total correction is at most $3\binom73$.  The selected
seven-circle contributes the remaining defect $27$.  Summing the block
inequality and using $C\leq72$ gives
\begin{align*}
 6W\geq{}&3\binom73+9\cdot7\binom72
 +12\binom72\cdot7+6\binom73+3\cdot21+6\cdot21\\
 &-36\cdot72-4\cdot140-27=412.
\end{align*}
Hence $W\geq69$.  On the other hand, for a fixed outsider pair the
$\Gamma$-pairs used by circles counted in $W$ are disjoint.  It contributes
at most $\floor{7/2}=3$, and therefore
\[
 W\leq3\binom72=63,
\]
a contradiction.
\end{proof}

\begin{theorem}\label{thm:large-range}
For every integer $n\geq14$,
\[
 f(n)=F(n)=1+\binom{n-1}{2}-\floor{\frac{n-1}{2}}.
\]
\end{theorem}

\begin{proof}
The lower bound for $n=14$ is \cref{prop:n14-dichotomy}.  For $n\geq15$,
\cref{lem:large-half-cap} supplies the common block cap, Langer supplies
\eqref{eq:large-L}, and the one functional in \cref{lem:large-SM} gives
the strict contradictions \eqref{eq:large-even-gap} and
\eqref{eq:large-odd-gap} under $C\leq F(n)-1$.  The foundational
circle-with-centre construction attains $F(n)$ for every $n\geq14$.
\end{proof}

\subsection{Fixed-certificate appendix and verification boundary}
\label{subsec:n13large-certificates}

For clarity, we record which parts of this section are fixed arithmetic
certificates and which parts are mathematical inputs.

\paragraph{The opening certificate.}
In \cref{lem:n13-gamma6-opening}, the residual
$123D_B-A(g,x)$ has the following line residuals:
\begin{align*}
 R_L(0,x)&=\frac{x(x-3)(890-47x)}2,\\
 R_L(1,x)&=\frac{(x-2)(-47x^2+597x+246)}2,\\
 R_L(2,x)&=\frac{x(x-1)(304-47x)}2.
\end{align*}
The circle residuals are
\begin{align*}
 R_C(0,x)&=237+\frac{(x-3)(-47x^2+644x-738)}2
                 -237\mathbf{1}_{\{x=3\}},\\
 R_C(1,x)&=\frac{(5-x)(47x^2-210x+390)}2
                 -237\mathbf{1}_{\{x=2\}},\\
 R_C(2,x)&=\frac{(x-2)(x-3)(284-47x)}2
                 -237\mathbf{1}_{\{x=1\}}.
\end{align*}
It is enough---and slightly stronger than the outer cap requires---to use
the intervals
\[
\begin{array}{c|ccc}
&g=0&g=1&g=2\\ \hline
\text{line}&3\leq x\leq7&2\leq x\leq6&1\leq x\leq5\\
\text{circle}&3\leq x\leq6&2\leq x\leq5&1\leq x\leq4.
\end{array}
\]
Each displayed expression is nonnegative on its indicated interval.  Thus
the opening is a coefficientwise proof, not a conclusion inferred from
sampled incidence patterns.

\paragraph{Affine certificate.}
For reference, retain the zero- and positive-slack variables
\[
 c_{03},c_{12},c_{15},c_{21},c_{22},l_{03},l_{12},l_{21}
\]
in addition to the variables in \cref{lem:n13-capacity-gap}.  With every
row written as left side minus right side, the eight conservation equations
are
\begin{align*}
R_C={}&1+c_{03}+h+k+A_6+c_{12}+y+z+c_{15}+c_{21}+c_{22}+f+g-(60-d),\\
R_O={}&c_{03}+c_{12}+c_{21}-(13+j),\\
R_0={}&c_{03}+4h+10k+20A_6+y+4z+10c_{15}+f+4g
       +l_{03}+4\ell+w+p-35,\\
R_1={}&c_{12}+3y+6z+10c_{15}+2c_{22}+6f+12g
       +l_{12}+3w+2t+6p-126,\\
R_2={}&c_{21}+2c_{22}+3f+4g+l_{21}+2t+3p-105,\\
R_W={}&c_{22}+3f+6g-(48+r),\\
R_X={}&3c_{03}-5k-12A_6+2c_{12}-4z-10c_{15}+c_{21}-3f-8g\\
 &\quad+3l_{03}+2l_{12}+l_{21}-3p-(21+u),\\
R_\Gamma={}&-12+c_{12}-z-2c_{15}+2c_{21}-2f-4g
       +l_{12}+2l_{21}-2p-(18+v).
\end{align*}
All eight equal zero.

\smallskip
\noindent\emph{The second-dual inequality.}
For a nonselected block $B$, put $g=|B\cap\Gamma|$, $x=|B\cap X|$ and
$s=g+x$.  Let $D_B=s(s-4)$ for a circle and $D_B=2s(s-2)$ for a line.
Weight the four triple rows by $6,9,12,6$, add $6M_X+3M_\Gamma$, and
subtract $36C+6W$.  The coefficient of $B$ is
\begin{align*}
 A_2(g,x)={}&6\binom x3+9g\binom x2+12\binom g2x
 +(6x+3g)(4-g-x)\\
 &-36\mathbf1_{\{B\text{ is a circle}\}}
 -6\mathbf1_{\{B\text{ is a circle and }g=2\}}\binom x2.
\end{align*}
For $g=0,1,2$, the residual $4D_B-A_2(g,x)$ is
\[
\begin{array}{c|ccc}
&g=0&g=1&g=2\\ \hline
\text{line}&
x(x-3)(14-x)&
(x-2)(-2x^2+21x+17)/2&
(x-1)(-x^2+7x+12)\\
\text{circle}&
(x-3)(-x^2+10x-12)&
(5-x)(x-2)(2x-3)/2&
(4-x)(x-2)(x-1).
\end{array}
\]
The line ranges are respectively $3\leq x\leq7$, $2\leq x\leq6$,
$1\leq x\leq5$, and the circle ranges are $3\leq x\leq6$,
$2\leq x\leq5$, $1\leq x\leq4$.  Every displayed factor is nonnegative;
the selected circle separately has equality $48=4\cdot12$.  Summing and
using $\mathsf D\leq117$ gives exactly
\begin{align}
 B_0:={}&12h+26k+36A_6+3y+5z+2f+40\ell+31w+22t+48p\notag\\
        &\quad+6u+3v+36d-(12+6r)\leq0.             \label{eq:n13-B0}
\end{align}
The zero-residual types are precisely
$c_{03},c_{12},c_{15},c_{21},c_{22},l_{03},l_{12},l_{21}$ and
$g=c(2,4)$; every omitted type has nonnegative residual and may safely be
discarded from the retained budget.
This proves the budget used in \eqref{eq:n13-second-budget}
coefficientwise.

\smallskip
\noindent\emph{The integral coordinate.}
In the row order $(R_C,R_O,R_0,R_1,R_2,R_W,R_X,R_\Gamma)$, define
\begin{align*}
 F_0={}&14f-\bigl(180A_6+24\ell-30r+36d+60h+36j+106k-140p\\
      &\qquad\quad-66t+2u-13v-21w+15y+z\bigr),\\
 G_0={}&336g-\bigl(294+342r-1716A_6-240\ell-276d-516h-276j\\
      &\qquad\quad-976k+1302p+618t-20u+123v+189w-87y+25z\bigr).
\end{align*}
Their fixed row certificates are
\begin{align}
 F_0={}&-36R_C+36R_O-6R_0+9R_1+24R_2-30R_W+2R_X-13R_\Gamma=0,
                                                               \label{eq:n13-F0}\\
 G_0={}&276R_C-276R_O+60R_0-83R_1-226R_2+342R_W
         -20R_X+123R_\Gamma=0.                                \label{eq:n13-G0}
\end{align}
Put $N=j+y+z$ and
\begin{align*}
 B_1={}&9N+36h+72k+108A_6+76\ell+49w+22t+49p\\
       &\quad+72d+11u+2v-(21+18r),\\
 H={}&2j+y-v+4h+7k+12A_6+\ell+2d-10p-5t-2w-2r.
\end{align*}
The right side subtracted from $14f$ in \eqref{eq:n13-F0} has the exact
decomposition
\[
 (K_0+7j-2r)+14H.
\]
Consequently
\[
 K_0+7j-2r=14(f-H).
\]
Set $\mu=f-H$ and $E_\mu=K_0+7j-2r-14\mu$.  Then $\mu$ is an integer,
$E_\mu=0$, and $\mu\geq0$: indeed the left side above is at least
$-2r\geq-12$, so it cannot be a negative multiple of $14$.

The remaining budget transformations are the fixed identities
\begin{align}
 7B_0-4B_1-F_0&=0,                                  \label{eq:n13-B01}\\
 B_1+21+18r&=2K_0+7Q_0,                             \label{eq:n13-B1KQ}\\
 B_1&=7\{Q_0+4\mu-2j-(3+2r)\}+2E_\mu.              \label{eq:n13-B1mu}
\end{align}
Thus $B_0\leq0$ and $F_0=E_\mu=0$ imply $B_1\leq0$ and precisely the two
relations in \eqref{eq:n13-mu}.

\smallskip
\noindent\emph{The $\delta$ and $J$ certificates.}
Let
\[
 \delta=7\mu-4j,\qquad
 L=3+2r-(Q_0+4\mu-2j)\geq0.
\]
Direct substitution in the displayed definitions gives the fixed identity
\begin{align}
 S&:=6-k-5\ell-6d-5t-5w-24p-21\mu+12j\notag\\
  &=f+2L+z\geq0,                                    \label{eq:n13-S}
\end{align}
and hence
\[
 3\delta=6-k-5\ell-6d-5t-5w-24p-S.
\]
This is \eqref{eq:n13-delta}.  The same substitution gives the two exact
equalities
\begin{align}
 v={}&2\delta+2r-a-2u-4h-8k-12A_6-10\ell-8d-4t-7w,
                                                               \label{eq:n13-v-cert}\\
 f={}&2j+\mu+y-v+4h+7k+12A_6+\ell+2d
       -10p-5t-2w-2r.                              \label{eq:n13-f-cert}
\end{align}

For the last elimination set $G=24g$ and
\begin{align*}
 G_*={}&21+42r+123\mu-228A_6-105\ell-90d-72h-83j-140k\\
      &\quad+93p+9t-19u-48w-8y-7N.
\end{align*}
Then the first fixed certificate is
\begin{equation}
 G_0-14(G-G_*)+123E_\mu=0,                          \label{eq:n13-Gstar}
\end{equation}
so $G=G_*$.  Let $E_\delta=\delta-(7\mu-4j)$, let $E_v$ be the left side
minus the right side of \eqref{eq:n13-v-cert}, and put
\begin{align*}
 E_{7f}=7f-\bigl(&18j+15\delta+7(y-a)-14v-7k-63\ell-42d\\
                 &-70p-63t-14u-63w\bigr).
\end{align*}
Equation \eqref{eq:n13-f-cert} gives $E_{7f}=0$.  Finally, let $E_J$ be
$7J$ minus the right side of \eqref{eq:n13-J}.  The terminal fixed
identity is
\begin{equation}
 E_J+7(G-G_*)-84E_v+7E_{7f}-123E_\delta=0.          \label{eq:n13-EJ}
\end{equation}
Every term except $E_J$ has just been shown to vanish, proving
\eqref{eq:n13-J}.  This explicit chain establishes
\eqref{eq:n13-mu}--\eqref{eq:n13-J}; no multiplier is chosen by
optimisation or incidence search.

\paragraph{What was checked arithmetically.}
During preparation, exact-arithmetic checks independently replayed the
opening residuals, the $W=47$
identities, all eight affine rows and their fixed eliminations, the seven
values of $\rho$, the capacity-gap conversion, the three-residue
calculations, the four half-cap polynomials, the coefficients of $S_M$, and
the fourteen-point residual table.  These checks manipulate integers, rationals,
and symbolic polynomials only.

Such arithmetic checks do not prove Melchior's theorem, Langer's inequality, the pivot form
of Lenchner's theorem~\cite{Lenchner2007}, the inversion dictionary, or the
real six-marked-conic signature and host lemmas.  Those are the stated
geometric inputs to the written proofs.  The classification-free
$\Gamma_5$ transfer in \cref{lem:n13-gamma5} is checked directly through
its displayed blockwise coefficients and endpoint equations.  The distributed
exact-arithmetic supplement
\path{supplement/verify_n13_gamma5_global_closure.py} checks an independent
global closure of that branch.  In every case the arithmetic is
corroborative, not a substitute for the theorem-level arguments above.
